\documentclass[11pt]{article}
\usepackage{amsmath,amssymb,amsthm}
\usepackage[margin=1in]{geometry}
\newtheorem{lemma}{Lemma}
\newtheorem{theorem}[lemma]{Theorem}
\newtheorem{corollary}[lemma]{Corollary}
\begin{document}
\title{A unified clique-extension lemma: proved research candidate}
\author{}\date{10 October 2026}\maketitle
This is an independent research candidate, not a modification of the main
manuscript or a formal verification. The arguments below are elementary
proofs. The separate finite experiment is not used to justify them.

Let $Y$ be a finite graph of order $M\ge3$, with finite labelled lists
$L_f$ for $f\in V(Y)$. Its clique extension $X$ has vertices $(f,c)$,
$c\in L_f$; each fibre is complete, and the cross edges are
$(f,c)(g,c)$ for $fg\in E(Y)$ and $c\in L_f\cap L_g$.
Assume $|L_f|\ge3$ and $|L_f\cap L_g|\ge2$ on every base edge.
For distinct labels $p,q$, let $B_{pq}$ be the spanning subgraph of $Y$
whose edges are exactly those with $L_f\cap L_g=\{p,q\}$.

\begin{lemma}[Exact obstruction for cyclic bridge selection]
Every base cycle admits bridge labels from its successive intersections,
with consecutive bridge labels distinct, if and only if every $B_{pq}$
is bipartite.
\end{lemma}
\begin{proof}
Lists of size at least two on a cycle admit a proper list colouring unless
the cycle is odd and every list is the same pair. For completeness, a list
of size at least three can be left until last; otherwise unequal consecutive binary
lists permit a greedy traversal beginning with a label absent from the
last list. Equal binary lists are coloured by alternation exactly on even
cycles. Thus a failing cycle is precisely an odd cycle of some $B_{pq}$.
\end{proof}

\begin{theorem}[Three clique-extension conclusions]
Suppose every $B_{pq}$ is bipartite. Consider also the following conditions:
\begin{enumerate}
\item[(E)] For every $f$ and $c\in L_f$, some neighbour $g$ of $f$ has
$c\in L_g$.
\item[(R)] Every base path containing all but at most one vertex has an
edge whose intersection list has size at least three.
\end{enumerate}
Then:
\begin{enumerate}
\item if $Y$ is Hamiltonian, $X$ is Hamiltonian;
\item if $Y$ is vertex-pancyclic and (E) holds, $X$ is vertex-pancyclic;
\item if $Y$ is Hamilton-connected and (E),(R) hold, $X$ is
Hamilton-connected. This third conclusion does not require the
bipartiteness assumptions.
\end{enumerate}
\end{theorem}
\begin{proof}
First, a path whose intersection lists have size at least two, with one
of size at least three, lifts with arbitrary prescribed endpoint labels. Choose
bridge labels greedily from the two ends towards the three-element list,
choosing there last. This avoids each prescribed endpoint and makes
successive bridge labels distinct. Spanning paths in the complete fibres
then concatenate.

For a cycle $f_0,\ldots,f_{s-1},f_0$, choose cyclic bridge labels by the
lemma. Spanning fibre paths give a spanning cycle on all its fibres,
proving conclusion 1 for a Hamilton base cycle. More generally, given
$x=(f_0,a)$, choose the path in its fibre on between $3$ and $|L_{f_0}|$
vertices, containing $x$, and the other fibre paths on between $2$ and
$|L_{f_i}|$ vertices. This gives every length in
\[
 [2s+1,\ \sum_{i=0}^{s-1}|L_{f_i}|]
\]
through $x$. Condition (E) gives an edge $f_0g$ preserving $a$ and
another common label $p\ne a$. The two cross edges at $a,p$, with
fibre paths on two or three vertices, give lengths $4,5,6$ through $x$;
its own fibre gives length $3$. For each $3\le s\le M$, take a base
cycle of length $s$ through $f_0$, and denote its total fibre size by
$T_s$. The intervals $[2s+1,T_s]$ overlap since $T_s\ge3s$ and
$2s+3\le3s$; moreover $T_M=|V(X)|$. This proves conclusion 2.

For conclusion 3, endpoints in different fibres are handled by a base
Hamilton path and (R). For endpoints $(f,a),(f,b)$ in one fibre, choose
$c\in L_f\setminus\{a,b\}$ and a neighbour $g$ preserving $c$ by (E).
Every edge of the Hamilton-connected base lies on a Hamilton cycle $C$.
Let $h$ be the other neighbour of $f$ on $C$ and choose
$q\in(L_f\cap L_h)\setminus\{c\}$. Interchange $a,b$ if $q=a$,
reversing the final path. Now split the root fibre into paths
$(f,a)\leadsto(f,c)$ and $(f,q)\leadsto(f,b)$; use a singleton for the
second when $q=b$, otherwise its endpoint edge, and put every remaining
vertex in the first. Condition (R) lifts $C-f$ through every other fibre
from $(g,c)$ to $(h,q)$. Join these three paths by the cross edges at
$c,q$.
\end{proof}

\begin{lemma}[Recolouring interfaces have the required bipartition]
Let $H=G-v\ne\varnothing$ and $k\ge\deg_G(v)+3$. With
$Y=\mathcal R_k(H)$ and $L_f=[k]\setminus f(N_G(v))$, every $B_{pq}$
is bipartite.
\end{lemma}
\begin{proof}
At an edge of $B_{pq}$, the one-coordinate change implies
$2\ge |L_f|-1$, so both endpoint lists have size three. Hence
$\deg_G(v)=k-3$, and its neighbours use pairwise distinct colours,
none equal to $p,q$. Order these neighbours. Append to their colour
tuple the unique unused colour in $[k]\setminus\{p,q\}$. This is a
permutation, and a $B_{pq}$ edge swaps one neighbour entry with the
last entry. Permutation sign gives a bipartition. Isolated vertices of
$B_{pq}$ may be assigned arbitrarily.
\end{proof}

\begin{corollary}[Application limits]
The preceding lemmas give one-vertex Hamiltonian lifting without any
global colour-permutation argument. They give vertex-pancyclic lifting
whenever (E) is available. They give Hamilton-connected lifting when
also $v$ has a non-neighbour and (E) is available.
\end{corollary}
\begin{proof}
A non-neighbour $w$ supplies (R): colour permutations give all $k\ge3$
labels at $w$, so any path omitting at most one base vertex must change
$w$, preserving the full available list. Condition (E) follows either
from $\delta(Y)>\deg_G(v)$ (at most $\deg_G(v)$ moves can destroy a
fixed available label), or from $H$ being nonempty and $d$-degenerate
with $k\ge d+3$ (recolour a low-degree vertex while avoiding its current
label and the fixed label).
\end{proof}

\paragraph{What remains.}
This unifies the existing lifting mechanism and isolates its exact
cyclic obstruction. It does not remove the complete-graph branch of
degenerate Hamilton-connectedness, because that branch has no
non-neighbour to guarantee (R). It does not establish that (R) may be
deleted. No new palette threshold is claimed.
\end{document}
